\documentclass[10pt,a4paper]{amsart}
\usepackage[margin=19mm]{geometry}
\usepackage{amsmath,amssymb,mathtools}
\usepackage[scr=rsfs,cal=euler]{mathalfa}
\usepackage{xfrac}
\usepackage[unicode,hidelinks,bookmarks=false]{hyperref}
\newcommand{\ie}{i.e.\ }
\newcommand{\cf}{cf.\ }
\newcommand{\resp}{respectively\ }
\newcommand{\Subsection}{\subsection}
\providecommand{\nobreakdash}{\nobreak\hbox{-}}
\setlength{\emergencystretch}{2em}
\multlinegap0pt
\usepackage{amssymb}
\usepackage{xcolor}
\newtheorem{prop}{Proposition}
\DeclareMathOperator{\End}{\mathrm{End}}
\newcommand{\R}{\mathbb R}
\def\co{\colon\thinspace}
\newcommand{\lra}{\longrightarrow}
\newcommand{\ra}{\rightarrow}
\newcommand{\rmd}{\mathrm d}
\title{A contact-geometric variant of Smale's contraction}
\author{Florian Buck, Christopher Schmidt, Kai Zehmisch}
\date{Source: arXiv:2609.06431v1 (CC BY 4.0)}
\begin{document}
\maketitle

\noindent\emph{Source and licence.} A contact-geometric variant of Smale's contraction. Version arXiv:2609.06431v1 (CC BY 4.0), distributed by arXiv under the Creative Commons Attribution 4.0 International licence, http://creativecommons.org/licenses/by/4.0/, as recorded in the arXiv licence metadata for this exact version and shown on the abstract page https://arxiv.org/abs/2609.06431v1. Full licence notice: \texttt{licenses/arxiv-2609.06431-CC-BY-4.0.txt}.\par\smallskip

\noindent\textit{Editorial scope.} Counted unit: the article's prop prop:globalflowbox (source0.tex lines 1032--1087). The article's complete original proof of it is delivered with it (source0.tex lines 1090--1189, one \texttt{proof} environment of 100 lines). No mathematics is written, repaired or re-derived: the statement and the proof are the article's own lines, in the article's own notation, and no retained line is substituted or re-flowed.  No further source-local context is needed: every cross-reference of every retained line resolves inside the two delivered spans.  Nothing was omitted from any retained span, so the package carries no omission record.  No retained line contains a citation, so the wrapper carries no bibliography and no bibliography entry.  No retained line contains a figure environment, an image-inclusion command or drawing code, so no image is delivered and the package declares no asset dependency.  Editorial formatting only, in addition: an AMS \texttt{amsart} preamble that restates the article's own declarations and macros used by the retained lines, namely the environments \texttt{prop} on separate counters, together with the standard packages those lines need.  The retained environments print in this document as Proposition 1 in order of appearance, whereas the article prints the same items with the numbering of its own full document.  The complete native source of the article as distributed in the arXiv e-print for this exact version is bundled unchanged as \texttt{sources/209534/source0.tex}.
\begin{prop}[Global Darboux]
 \label{prop:globalflowbox}
 The following statements are equivalent:
	\begin{enumerate}
	  \item[(a)]
	  $(M,\alpha)$ is vertically convex.
	  \item[(b)]
	  $(\hat{M},\hat{\alpha})$ has no Reeb orbits
	  trapped in backward time.
	  \item[(c)]
	  There exist a shape $(T,V_T,g_{\pm})$
	  equivalent to $(S,V_S,f_{\pm})$ with $g_-=f_-$
	  and an intrinsically defined strict contactomorphism
	  \[
	    \Phi\co
	    \big(\R\times V,D_T,\rmd b+\lambda\big)
	    \lra
	    (\hat{M},M,\hat{\alpha})
	    \,,
	  \]
	  that sends $\big(\!\End(D_T),\rmd b+\lambda\big)$
	  to $\big(\!\End(D_S),\rmd b+\lambda\big)$.
	  The restriction to
	  \[
	    Z_+:=\{b\geq g_+\}\subset\R\times V_T
	  \]
	  is of the form
	  \[
	    \Phi(b,v)=\big(b+f(v),\varphi(v)\big)
	  \]
	  for all $(b,v)\in Z_+$.
	  Here $\varphi$ is a compactly supported,
	  exact symplectomorphism of $(V_T,\rmd\lambda)$
	  with
	  \[
	  \varphi^*\lambda=\lambda-\rmd f
	  \]
	  and the function
	  \[
	  f=\varphi^*f_+-g_+
	  =(\varphi^*f_+-f_-)-(g_+-g_-)
	  \]
	  has compact support in $V_T$ and determines $g_+$
	  uniquely.
	  The restriction of $\Phi$ to $\End(D_T)\setminus Z_+$
	  is the identity map.
	  
	  In fact,
	  $\Phi(b,v)=\big(b+f(v),\varphi(v)\big)$
	  holds for all $(b,v)$ in the neighbourhood of $\End(D_T)$
	  determined by $\Phi^{-1}\big(\psi(U)\big)$,
	  where $f(v)=0$ and $\varphi(v)=v$
	  whenever $(b,v)$ lies
	  in or sufficiently close to $\End(D_T)\setminus Z_+$.
	\end{enumerate}
\end{prop}

\begin{proof}
  Because $(D_T,\rmd b+\lambda)$
  has no trapped Reeb orbits,
  the implication (c) $\Rightarrow$ (a) follows.
  The implication (a) $\Rightarrow$ (b)
  follows directly from the definition. 
  
  We are left with (b) $\Rightarrow$ (c).
  Denote by $\varphi_t$ the Reeb flow of $\hat{\alpha}$,
  which is complete by construction.
  Denote by $b_-$ the minimum of $f_-$.
  Then $\{b_-\}\times V$ is a well-defined
  hypersurface in $\hat{M}$ transverse
  to the Reeb flow $\varphi_t$.
  Using $\{b_-\}\times V$ as a Cauchy hypersurface,
  we define a map $\Phi\co\R\times V\ra\hat{M}$ by setting
  \[
    \Phi(b,v):=\varphi_{b-b_-}(b_-,v)
    \,.
  \]
  The map is injective,
  because each Reeb orbit intersects
  the separating hypersurface $\{b_-\}\times V$ at most once,
  as all intersections are transverse and
  occur with the same co-orientation.
  The map is surjective,
  because all Reeb orbits intersect $\{b_-\}\times V$
  by assumption (b).
  Therefore, $\Phi^{-1}$ exists.
  The differential $T_{(b,v)}\Phi$ evaluates on $\partial_b$ to
  the Reeb vector field $R_{\hat{\alpha}}$ of $\hat{\alpha}$
  taken at $\varphi_{b-b_-}(b_-,v)$;
  it evaluates on $T(\{b\}\times V)$ to $T_v\varphi_{b-b_-}(b_-,\,.\,)$. 
  Since $\varphi_t$ is a diffeomorphism for each $t$,
  the Reeb vector field $R_{\hat{\alpha}}$ is transverse
  to the level sets $\varphi_{b-b_-}(b_-,V)$.
  Hence, $\Phi$ is an immersion.
  The implicit function theorem yields smoothness
  of the inverse $\Phi^{-1}$.
  Moreover,
  $\Phi^*\hat{\alpha}$ evaluates on $\partial_b$ to $1$
  as we integrate $R_{\hat{\alpha}}$;
  it restricts to $\lambda$ on $T(\{b\}\times V)$ for all $b\in\R$
  because the Reeb flow of $\hat{\alpha}$ preserves $\hat{\alpha}$
  and the restriction of $\hat{\alpha}$ to $T(\{b_-\}\times V)$
  equals $\lambda$.
  Hence, $\Phi^*\hat{\alpha}=\rmd b+\lambda$.
    
  We verify the remaining properties of $\Phi$ under assumption (b).
  The hypersurface $T:=\Phi^{-1}(S)$
  in $\big(\R\times V,\rmd b+\lambda\big)$
  is transverse to $\partial_b$ along $T\setminus\Phi^{-1}(E_S)$.
  Since $\Phi$ is a diffeomorphism and $M$ is compact,
  the domain
  \[
  D_T:=\Phi^{-1}(M)
  \]
  is compact.
  By construction,
  $\Phi(b,v)=(b,v)$ for all $(b,v)\in\R\times V_S$
  with $b\leq f_-(v)$.
  Hence, for every $v\in V_S$,
  the vertical line through $v$ enters $D_T$
  through $\big(f_-(v),v\big)$
  and can leave it only through $\Phi^{-1}(S^+)$.
  Its intersection with $D_T$ is therefore a compact interval.
  By the implicit function theorem,
  its upper endpoint depends smoothly on $v$.
  Thus, $T$ is a strict vertically convex hypersurface
  with shadow set $V_T=V_S$ and lower graph $g_-=f_-$.
  The shapes $(T,V_T,g_{\pm})$ and $(S,V_S,f_{\pm})$
  are equivalent with $g_-=f_-$,
  because $T$ coincides with $S$
  along a neighbourhood of the closure of $S^-$.
  
  By construction,
  $\Phi$ is the identity map
  in a neighbourhood of $\End(D_T)\setminus Z_+$.
  On a neighbourhood of $Z_+$,
  the map $\Phi$ is a strict contactomorphism
  onto its image with respect to $\rmd b+\lambda$.
  In particular, $\Phi_*\partial_b=\partial_b$.
  Hence,
  $\Phi(b,v)=\big(b+f(v),\varphi(v)\big)$
  for all $(b,v)$ in this neighbourhood.
  Notice that $f$ has compact support in $V_T$.
  A repetition of the argument with $\Phi^{-1}|_{\End(D_S)}$
  shows that $\varphi$ is a compactly supported
  diffeomorphism of $V_T$.
  Moreover, since $\rmd b+\lambda$ equals
  $\Phi^*\big(\rmd b+\lambda\big)=\rmd b+\rmd f+\varphi^*\lambda$
  on a neighbourhood of $Z_+$,
  we get $\varphi^*\lambda=\lambda-\rmd f$ on $V_T$.
  Finally,
  for all $v\in V_T$ there exists a unique $w\in V_S$
  such that $\Phi\big(g_+(v),v\big)=\big(g_+(v)+f(v),\varphi(v)\big)$
  is equal to $\big(f_+(w),w\big)$.
  Hence,
  $w=\varphi(v)$ and $\varphi^*f_+=g_++f$.
\end{proof}

\end{document}
